% =====================================================================
% Section: Introduction (with Related Work)  (outline.md §1)
% Draft v1 (2026-06-10). Related Work is integrated into the Introduction
% following the CS/engineering SCI body convention (outline.md).
% Citation keys refer to paper/references.bib.
% This draft contains NO quantitative result assertions; all comparative
% statements are qualitative (motivation / positioning) and are backed by
% the experiments referenced in plan/review/method-experiment-traceability.md.
% =====================================================================
\section{Introduction}
\label{sec:intro}

Single image super-resolution (SR) aims to reconstruct a high-resolution (HR)
image from its low-resolution (LR) observation, and it underpins a wide range of
applications from mobile photography to display rendering. Since the seminal
deep model SRCNN~\cite{dong2014srcnn}, accuracy has improved steadily with
deeper and wider networks such as EDSR~\cite{lim2017edsr}. However, the
accompanying growth in parameters and computation makes these models impractical
for resource-constrained platforms such as mobile and edge devices, where memory,
latency, and energy are tightly budgeted. This has motivated a sustained line of
work on \emph{lightweight} SR, which seeks a favorable trade-off between
reconstruction quality and model cost (parameters, FLOPs, and inference latency).

\paragraph{Lightweight CNN-based SR.}
Early efficient designs reduce cost through architectural economy. CARN employs
cascading residual connections with group convolutions for a fast and compact
network~\cite{ahn2018carn}. IMDN introduces information multi-distillation to
progressively refine and aggregate features under a tight budget~\cite{hui2019imdn},
and RFDN reformulates this idea with residual feature distillation
blocks~\cite{liu2020rfdn}. RLFN further simplifies the distillation topology to
lower latency while preserving accuracy, winning the NTIRE efficient-SR
challenge~\cite{kong2022rlfn}. These methods are highly efficient, but their
purely convolutional, locally-connected operators have a limited receptive field
and cannot easily model the long-range, repetitive structure that natural and
graphical images exhibit.

\paragraph{Lightweight Transformer-based SR.}
Self-attention captures long-range dependencies and has been adapted to efficient
SR. SwinIR introduces window-based attention for image
restoration~\cite{liang2021swinir}; ELAN streamlines long-range attention to cut
redundant computation~\cite{zhang2022elan}; and SRFormer revisits the
window--channel trade-off with permuted self-attention~\cite{zhou2023srformer}.
ATD augments local windows with an adaptive token dictionary to extend the
effective context~\cite{zhang2024atd}. While effective, window attention couples
the receptive field to a fixed spatial partition, so tokens that are semantically
similar but spatially distant are not directly compared unless the window grows,
which raises cost.

\paragraph{Content-aware / cluster-based routing.}
A complementary direction groups tokens by \emph{content} rather than by spatial
position, so that long-range but semantically related regions interact within a
compact attention. SPIN performs superpixel-based token interaction for
lightweight SR~\cite{zhang2024spin}. CATANet pushes this idea further with a Token
Aggregation Block (TAB) that maintains a small set of routing centers, assigns
each token to its nearest center, reorders tokens by cluster label, and aggregates
the reordered sequence with an intra/inter-group attention
operator~\cite{liu2025catanet}. This content routing delivers strong accuracy at a
low budget and is the most direct predecessor of our work.

\paragraph{Limitations of center-based routing.}
Despite its effectiveness, the TAB routing mechanism has four limitations.
(i) Its routing centers are maintained across batches through a $k$-means--style
iteration with an exponential moving average (EMA) buffer, so the centers are
driven by a cross-batch historical state rather than strictly by the current
image, and adapt only slowly to per-image semantics. (ii) The token-to-center
relation is a hard \texttt{argmax} label that discards the \emph{confidence} of
the assignment, treating a token that clearly belongs to one cluster identically
to an ambiguous one on a boundary. (iii) The reordering clusters tokens only by
label, with no intra-cluster purity constraint, so high- and low-confidence tokens
are interleaved within a group during aggregation. (iv) Prototype \emph{generation}
and \emph{confirmation} are entangled in a single iterative update, which weakens
interpretability and stability. These observations suggest that a routing module
built directly from the current input, that preserves and exploits assignment
confidence, and that decouples prototype generation from confirmation, could
improve content routing without sacrificing the lightweight regime.

\paragraph{Our approach.}
We propose \emph{Dynamic Prototype Routing} (DPR), a drop-in replacement for the
TAB routing that keeps the CATANet backbone, its external tensor interface, and
its lightweight parameter budget unchanged. DPR generates a set of semantic
prototypes directly from the current tokens via a soft assignment, confirms them
with a learnable query bank decoupled from generation, matches tokens back to the
confirmed prototypes with a learnable temperature, and propagates an explicit
per-token confidence through ordering and aggregation. The resulting network,
\textbf{DPRNet}, is trained on DIV2K and evaluated on five standard benchmarks at
$\times2/\times3/\times4$, where it matches or slightly surpasses the strongest
content-routing baseline across benchmarks at a comparable parameter and
computation budget, while exposing an interpretable, confidence-aware routing
mechanism.

\paragraph{Contributions.}
Our main contributions are summarized as follows.
\begin{itemize}
  \item \textbf{C1 -- Dynamic semantic prototypes.} We generate routing prototypes
  from the current tokens by a soft, assignment-weighted aggregation, removing the
  cross-batch EMA center buffer so that prototypes follow per-image semantics
  (Sec.~\ref{sec:method_dpr}, Eqs.~\eqref{eq:assign}--\eqref{eq:proto_content}).
  \item \textbf{C2 -- Decoupled prototype confirmation.} We separate prototype
  \emph{generation} from \emph{confirmation} by refining the content prototypes
  with a learnable query bank through a single gated cross-attention, improving the
  stability and interpretability of the routing slots
  (Eq.~\eqref{eq:proto_refine}).
  \item \textbf{C3 -- Confidence-aware routing.} We explicitly retain the token
  assignment confidence and use it in three places: a composite ordering key that
  sorts high-confidence tokens first within each cluster, a confidence-modulated
  global gate, and a soft prototype fallback for uncertain tokens
  (Eqs.~\eqref{eq:sortkey}, \eqref{eq:iasa_gate}, \eqref{eq:soft_fallback}).
  \item \textbf{C4 -- Stabilized routing.} We introduce a learnable routing
  temperature together with a prototype usage-balancing regularizer to prevent
  collapse onto a few prototypes and to reduce seed-to-seed variance
  (Eqs.~\eqref{eq:scores}, \eqref{eq:balance}).
\end{itemize}
Each design element is exposed as an independent switch that defaults to the full
model and is exactly equivalent to its disabled form when turned off, so the
contributions can be isolated cleanly in the ablation study
(Sec.~\ref{sec:ablation}).
